We present the complete formal mathematical model underlying SheafDB.
This section consolidates notation, definitions, and structural properties.
All definitions are standard within sheaf theory~\cite{maclane1992sheaves}
and are adapted to the finite, discrete setting of knowledge representation.

\subsection{Formal Notation}
\label{sec:notation}

We use the following notation throughout:

\begin{itemize}
    \item $\mathbb{N}$: natural numbers $\{0,1,2,\ldots\}$.
    \item $\mathcal{P}(X)$: power set of $X$.
    \item $|X|$: cardinality of $X$.
    \item $f: A \to B$: total function; $f: A \rightharpoonup B$: partial function.
    \item $g \circ f$: composition $(g \circ f)(x) = g(f(x))$.
    \item $\operatorname{id}_X$: identity function on $X$.
    \item $(P, \leq)$: a poset with order relation $\leq$.
    \item $\downarrow x = \{y \in P \mid y \leq x\}$: down-set of $x$.
    \item $x \wedge y$: meet (greatest lower bound); $x \vee y$: join (least upper bound).
    \item $\mathcal{I}$: set of identifiers (UUIDs).
    \item $\mathcal{E} \subseteq \mathcal{I}$: set of entity identifiers.
    \item $\mathcal{R} \subseteq \mathcal{I}$: set of relation identifiers.
    \item $\mathcal{V} = \mathcal{I} \sqcup \mathbb{L}$: values (identifiers or literals).
    \item $\mathcal{C}$: set of contexts, a finite poset.
    \item $\mathcal{F}$: set of all possible semantic facts.
    \item $F$: a presheaf or sheaf.
    \item $F(U)$: sections over open set $U$.
    \item $\rho_{V,U}: F(V) \to F(U)$: restriction map for $U \subseteq V$.
    \item $F_x$: stalk at point $x$: $\varinjlim_{U \ni x} F(U)$.
    \item $\Gamma(X, F)$: global sections $F(X)$.
    \item $\llbracket Q \rrbracket$: denotation (result) of query $Q$.
\end{itemize}

\subsection{Category-Theoretic Setting}
\label{sec:categorical}

Let $\mathbf{Set}$ denote the category of sets and functions,
and $\mathbf{Pos}$ the category of partially ordered sets and monotone maps.
A presheaf on a small category $\mathcal{C}$ is a functor
$F: \mathcal{C}^{\mathrm{op}} \to \mathbf{Set}$.

In our setting, $\mathcal{C}$ is the context poset $(\mathcal{C}, \leq)$
viewed as a category: objects are contexts, and there is a unique morphism
$c \to d$ iff $d \leq c$. The opposite category $\mathcal{C}^{\mathrm{op}}$
reverses this: a morphism corresponds to $c \leq d$, i.e., restriction from
a broader context to a narrower one.

\subsection{Formal Definitions}
\label{sec:formal_definitions}

\begin{definition}[SemanticFact]
\label{def:semanticfact}
A \emph{semantic fact} is a tuple
\[
f = (i, s, r, \vec{o}, c, k, m)
\]
where:
\begin{itemize}
    \item $i \in \mathcal{I}$ is a globally unique identifier.
    \item $s \in \mathcal{E}$ is the subject entity.
    \item $r \in \mathcal{R}$ is the relation type.
    \item $\vec{o} \in \mathcal{V}^*$ is an ordered sequence of values (arity $\operatorname{arity}(f) = |\vec{o}|$).
    \item $c \in \mathcal{C}$ is the context of validity.
    \item $k \in [0,1]$ is the confidence.
    \item $m \in \mathbf{Set}$ is metadata.
\end{itemize}
Let $\mathcal{F}$ denote the set of all semantic facts.
\end{definition}

$\mathcal{E} \subseteq \mathcal{I}$ and $\mathcal{R} \subseteq \mathcal{I}$
denote the sets of entities and relations respectively, each carrying a
type and (for relations) an arity that a fact's object tuple must match;
Appendix~\ref{sec:appendix:definitions} gives the full structure, which is
not needed to follow the rest of the paper.

\begin{definition}[Context]
\label{def:context}
A \emph{context} $c \in \mathcal{C}$ is a node in a finite poset $(\mathcal{C}, \leq)$
identified by a dot-separated path. The partial order is prefix inclusion:
$c_1 \leq c_2$ iff $c_2$ is a prefix of $c_1$ (i.e., $c_1$ is more specific).
The root context $\top$ (``world'') is the unique maximal element:
$c \leq \top$ for all $c \in \mathcal{C}$.
Joins always exist: $c_1 \vee c_2$ (the least upper bound) is the longest
common prefix, i.e.\ the most specific common ancestor. Meets are partial:
if $c_1 \leq c_2$ then $c_1 \wedge c_2 = c_1$, but two incomparable
contexts have no common lower bound at all, since no context can carry two
distinct sibling paths as prefixes. The context hierarchy is therefore a
tree-shaped join-semilattice, not a lattice
(Appendix~\ref{sec:appendix:theorems} states the resulting open-set
properties).
\end{definition}

\begin{definition}[Finite Topological Space]
\label{def:topology}
A \emph{finite topological space} is a pair $(X, \mathcal{T})$ where $X$ is a finite set
and $\mathcal{T} \subseteq \mathcal{P}(X)$ satisfies:
\begin{enumerate}
    \item $\emptyset, X \in \mathcal{T}$.
    \item $U, V \in \mathcal{T} \implies U \cap V \in \mathcal{T}$ (finite intersection).
    \item $\{U_i\}_{i \in I} \subseteq \mathcal{T} \implies \bigcup_i U_i \in \mathcal{T}$ (arbitrary union).
\end{enumerate}
In SFDB, $X = \mathcal{I}$ (fact identifiers) and $\mathcal{T}$ is generated
by the Alexandrov topology on $(\mathcal{C}, \leq)$.
\end{definition}

\paragraph{Remark.}
Section~\ref{sec:sheaf-remark}, immediately below, explains why the sheaf
condition holds vacuously on this topology and is not the paper's
technical contribution; we do not repeat that argument here.

An \emph{open set} $U \in \mathcal{T}$ is a subset of $X$ declared open; in
the Alexandrov topology induced by $(\mathcal{C}, \leq)$, every down-set
$\downarrow c = \{d \in \mathcal{C} \mid d \leq c\}$ is open, and each
context $c$ has an associated open set
$U_c = \{f \in \mathcal{I} \mid \operatorname{context}(f) \leq c\}$
(Appendix~\ref{sec:appendix:definitions} gives this formally, together
with its intersection and top-element properties).

\begin{definition}[Neighbourhood]
\label{def:neighbourhood}
A \emph{neighbourhood} of $x \in X$ is any $U \in \mathcal{T}$ with $x \in U$.
The neighbourhood filter is $\mathcal{N}(x) = \{U \in \mathcal{T} \mid x \in U\}$.
In the Alexandrov topology, the minimal neighbourhood exists uniquely:
$\mathcal{B}(x) = \bigcap_{U \in \mathcal{N}(x)} U = \downarrow \operatorname{context}(x)$,
using the same context-assignment function $\operatorname{context}: \mathcal{I} \to \mathcal{C}$
that defines $U_c$ above.
\end{definition}

\begin{definition}[Restriction Map]
\label{def:restriction}
For $c \leq d$ (so $U_c \subseteq U_d$), the \emph{restriction map}
$\rho_{d,c}: F(d) \to F(c)$ is the identity on facts, restricted to the
subset of $F(d)$ that already lies in $U_c$:
\[
\rho_{d,c}(f) =
\begin{cases}
f & \text{if } \operatorname{context}(f) \leq c, \\
\text{undefined} & \text{otherwise.}
\end{cases}
\]
A fact is not rewritten by restriction --- no field of $f$ changes --- it
is simply re-indexed as an element of the smaller set $F(c)$ when its own
context already places it there. This is the presheaf of \emph{visibility
from an open set}: restricting a broader view to a narrower one either
finds the fact was already visible from that narrower vantage point, or
it was not, in which case restriction is undefined rather than
approximate. (An earlier version of this definition described
restriction as dropping object slots not meaningful in the target
context; the implementation does not perform any such transformation,
and that description is retracted here rather than left standing
alongside code that does something else.)
Functoriality: $\rho_{c,c} = \operatorname{id}_{F(c)}$, and for
$c \leq d \leq e$, $\rho_{e,c} = \rho_{d,c} \circ \rho_{e,d}$ holds
because both sides equal the identity restricted to
$\{f \in F(e) : \operatorname{context}(f) \leq c\}$.
\end{definition}

\begin{definition}[Stalk]
\label{def:stalk}
The \emph{stalk} of presheaf $F$ at point $x \in X$ is
$F_x = \varinjlim_{U \in \mathcal{N}(x)} F(U)$,
the set of germs $[s]_U$ where $[s]_U \sim [t]_V$ iff $\exists W \subseteq U \cap V$
with $x \in W$ and $s|_W = t|_W$.
\end{definition}

On an Alexandrov space this direct limit is not an abstract object: every
point has a minimal open neighbourhood $\mathcal{B}(x)$
(Definition~\ref{def:neighbourhood}), the limit collapses to
$F_x \cong F(\mathcal{B}(x))$, and no germ equivalence needs to be
computed. This is what licenses the implementation's \emph{stalk index}:
a hash table mapping each fact identifier $x$ to the section over its
minimal neighbourhood is literally a table of stalks, materialised in
advance, and the paper's ``stalk lookup'' is a lookup in that table
rather than a direct-limit computation at query time.

\begin{definition}[Local Section]
\label{def:localsection}
A \emph{local section} is $s \in F(U)$ for a proper open subset $U \subsetneq X$.
In SFDB, a local section is a fact whose context $c < \top$.
\end{definition}

\begin{definition}[Global Section]
\label{def:globalsection}
A \emph{global section} is $s \in F(X)$ over the whole space.
In SFDB, this is a fact valid in the root context $\top$.
Computed by collecting local sections, checking compatibility on overlaps,
and gluing upward until $\top$ is reached.
\end{definition}

A \emph{knowledge state} $\Sigma = (X, F)$ pairs a finite topological space
of fact identifiers with a presheaf $F$ assigning facts to open sets, with
insert/delete/update as set operations on $\Sigma$; a \emph{query}
$Q = (\tau, s, r, \vec{o}, c, k_{\min}, n_{\max})$ denotes the set of facts
$\llbracket Q \rrbracket_\Sigma$ matching its subject/relation/object/context/confidence
pattern. Both are given fully in Appendix~\ref{sec:appendix:definitions}.
Section~\ref{sec:correctness} defines the notion of \emph{result
equivalence} actually used by the implementation and the equivalence
proof (Definition~\ref{def:semanticequiv}); it is the only such
definition in this paper.

\subsection{Theorems}
\label{sec:theorems}

\subsubsection{Semantic Preservation}
\label{sec:thm:preservation}

\begin{theorem}[Semantic Preservation / Round-trip Correctness]
\label{thm:preservation}
Let $T: \mathcal{F} \to \mathcal{T}(\mathcal{F})$ be the triple decomposition
function (reifying an $n$-ary fact into $2+n$ triples), and
$R: \mathcal{T}(\mathcal{F}) \to \mathcal{F}$ be the reconstruction function.
For any semantic fact $f \in \mathcal{F}$:
\[
R(T(f)) = f.
\]
\end{theorem}

\begin{proof}[Proof Sketch]
Decomposition $T$ creates triples $(s, \mathtt{rdf:type}, \mathtt{rdf:Statement})$,
$(s, \mathtt{rdf:subject}, f.\text{subject})$, $(s, \mathtt{rdf:predicate}, f.\text{relation})$,
and $(s, \mathtt{rdf:object}_j, f.\text{objects}_j)$ for each $j$.
Reconstruction $R$ groups triples by statement ID $s$ and inverts each step.
Since $T$ is injective and $R$ is its left inverse on $\operatorname{im}(T)$,
the round-trip is identity.
\end{proof}

\subsubsection{Functoriality of Restriction}

\begin{theorem}[Restriction Functoriality]
\label{thm:restriction}
The restriction maps $\rho_{d,c}: F(d) \to F(c)$ for $c \leq d$ satisfy:
\[
\rho_{c,c} = \operatorname{id}_{F(c)}, \qquad
\rho_{e,c} = \rho_{d,c} \circ \rho_{e,d} \quad (c \leq d \leq e).
\]
Hence $F: \mathcal{C}^{\mathrm{op}} \to \mathbf{Set}$ is a functor (a presheaf).
\end{theorem}

\subsubsection{Sheaf Axioms}

\begin{theorem}[Locality]
\label{thm:locality}
Let $F$ be a presheaf on $(X, \mathcal{T})$, $U \in \mathcal{T}$,
$\{U_i\}$ an open cover of $U$, and $s, t \in F(U)$.
If $\rho_{U,U_i}(s) = \rho_{U,U_i}(t)$ for all $i$, then $s = t$.
\end{theorem}

\begin{theorem}[Gluing]
\label{thm:gluing}
Let $F$ be a sheaf on $(X, \mathcal{T})$, $U \in \mathcal{T}$,
$\{U_i\}$ an open cover of $U$, and $\{s_i \in F(U_i)\}$ a compatible family:
$\rho_{U_i, U_i \cap U_j}(s_i) = \rho_{U_j, U_i \cap U_j}(s_j)$ for all $i,j$.
Then there exists a unique $s \in F(U)$ with $\rho_{U,U_i}(s) = s_i$ for all $i$.
\end{theorem}

\begin{proof}[Proof Sketch]
For each $x \in U$, pick $i$ with $x \in U_i$ and define $s(x) = s_i(x)$.
Compatibility ensures this is well-defined on overlaps.
Uniqueness follows from the locality axiom.
\end{proof}

\subsubsection{Canonical Equivalence}

\begin{theorem}[Invertibility of the Canonical Mapping]
\label{thm:canonical}
Let $\phi: \mathcal{F} \to \mathfrak{F}$ be the canonical mapping from facts to canonical facts,
and $\psi: \mathfrak{F} \to \mathcal{F}$ the inverse mapping. Then:
\[
\psi(\phi(f)) = f,\quad \phi(\psi(\mathfrak{f})) = \mathfrak{f}
\]
for all $f \in \mathcal{F}$, $\mathfrak{f} \in \mathfrak{F}$.
\end{theorem}

\begin{theorem}[Cross-Engine Query Equivalence]
\label{thm:equivalence}
Let $\Sigma_{\mathrm{KG}}$ and $\Sigma_{\mathrm{Sheaf}}$ be knowledge states
derived from the same canonical model $\mathfrak{F}$.
Let $Q$ be a canonical query. Then:
\[
\llbracket Q \rrbracket_{\Sigma_{\mathrm{KG}}} \cong \llbracket Q \rrbracket_{\Sigma_{\mathrm{Sheaf}}}
\]
where $\cong$ denotes equality up to identifier isomorphism.
\end{theorem}

\subsubsection{Structural Properties}

\begin{theorem}[Context Meet as Open Set Intersection]
\label{thm:meet}
For any $c_1, c_2 \in \mathcal{C}$ in the Alexandrov topology $\mathcal{T}_A$:
\[
U_{c_1} \cap U_{c_2} = U_{c_1 \wedge c_2}.
\]
\end{theorem}

\begin{proof}
$d \in U_{c_1} \cap U_{c_2}$ iff $d \leq c_1$ and $d \leq c_2$, which is
equivalent to $d \leq c_1 \wedge c_2$, i.e., $d \in U_{c_1 \wedge c_2}$.
\end{proof}

\begin{theorem}[Preservation of Referential Integrity]
\label{thm:integrity}
For any knowledge state $\Sigma$ and fact $f \in \Sigma$:
\[
\operatorname{subject}(f) \in \mathcal{E}_\Sigma,\quad
\operatorname{relation}(f) \in \mathcal{R}_\Sigma,\quad
\forall \omega_j \in \vec{o}: \omega_j \in \mathcal{E}_\Sigma \text{ if reference}.
\]
\end{theorem}

\begin{theorem}[Determinism of Query Execution]
\label{thm:determinism}
For a fixed query $Q$ and fixed knowledge state $\Sigma$,
the result $\llbracket Q \rrbracket_\Sigma$ is uniquely determined,
independent of execution order, caching state, parallelism, or query history.
\end{theorem}

\begin{theorem}[Consistency of Global Sections]
\label{thm:global}
A family of local sections $\{s_i \in F(U_i)\}$ can be glued to a global section
$s \in F(\mathcal{C})$ iff for every pair $(i,j)$:
\[
\rho_{U_i, U_i \cap U_j}(s_i) = \rho_{U_j, U_i \cap U_j}(s_j).
\]
\end{theorem}




